% Thesis v3 (viet tu thesis loi V3, 2026-11-11). Moi muc chi commit sau khi Le duyet.
% Ch1 theo ban Le viet 2026-11-12 (1.1, 1.2, 1.3, 1.5); 1.4 va 1.6 giu tu ban truoc.
\chapter{Introduction}\label{ch:introduction}
\definecolor{passc}{RGB}{33,102,172}
\definecolor{failc}{RGB}{196,78,52}
\definecolor{focusc}{RGB}{33,102,172}
\tikzset{
  ctag/.style={circle, draw=black!60, fill=white, inner sep=0.6pt, font=\tiny\bfseries},
}
\definecolor{accentc}{RGB}{214,150,0}


\section{Overview}\label{sec:intro-background}

A 2024 analysis by Google estimates that memory-safety bugs account for approximately 70\% of severe vulnerabilities in large C and C++ codebases~\cite{rebert2024memorysafety}. Against this backdrop, the Defense Advanced Research Projects Agency (DARPA) launched the Translating All C to Rust (TRACTOR) program in 2024 to automate the migration of legacy C software to Rust, aiming to address memory-safety vulnerabilities while preserving program functionality~\cite{darpa-tractor}. Rust offers a promising migration target because its ownership and borrowing mechanisms prevent many such defects in safe code~\cite{rustreference}. However, migrating existing C programs remains challenging, as differences in memory management, pointer semantics, and type systems must be reconciled without altering the behavior on which existing applications depend.

Recent advances in large language models (LLMs) have opened new possibilities for automating C-to-Rust migration. LLM-based translation systems can generate Rust code from C and iteratively refine it using compilation errors, program analysis, and execution feedback~\cite{nitin2025c2saferrust,shiraishi2026smartc2rust,schesch2026actor}.

\begin{figure}[htbp]
\centering
\resizebox{\linewidth}{!}{\begin{tikzpicture}[>=Stealth, font=\footnotesize,
  st/.style={draw, rounded corners=2pt, align=center, minimum height=10mm, inner sep=4pt, fill=white},
  io/.style={st, fill=black!6}]
% ---- translation, generalized (not studied) ----
\node[io, minimum width=2.4cm] (c) at (0.6,0) {C program};
\node[st, minimum width=5.0cm, minimum height=13mm, fill=black!3, dashed] (tr) at (5.0,0)
  {\textbf{Translator}\\[1pt]\scriptsize produces Rust code; its internal\\[-1pt]\scriptsize steps differ from system to system};
\node[io, minimum width=2.6cm] (rs) at (9.6,0) {Rust translation};
\draw[->, thick] (c) -- (tr); \draw[->, thick] (tr) -- (rs);
% ---- evaluation, expanded ----
\node[st, minimum width=3.6cm] (cmp) at (9.6,-1.9) {Compile the translation};
\node[st, fill=black!12, text=black!70, minimum width=2.6cm] (fb) at (14.7,-1.9) {FAIL at the\\build stage};
\node[io, minimum width=3.2cm] (tests) at (4.3,-3.55) {Benchmark tests:\\inputs and expected results};
\node[st, minimum width=3.6cm] (run) at (9.6,-3.55) {Run on each test input;\\compare with the expected result};
\node[st, very thick, draw=passc, fill=passc!10, minimum width=3.0cm] (pass) at (7.6,-5.3) {\textbf{\textcolor{passc}{PASS}}\\every result matches};
\node[st, very thick, draw=failc, fill=failc!10, minimum width=3.0cm] (ft) at (11.6,-5.3) {\textbf{\textcolor{failc}{FAIL}} at the test stage\\some result differs};
\draw[->, thick] (rs) -- (cmp);
\draw[->, thick] (cmp) -- node[above, align=center, font=\scriptsize]{does not\\compile} (fb);
\draw[->, thick] (cmp) -- node[right, font=\scriptsize]{compiles} (run);
\draw[->, thick] (tests) -- (run);
\draw[->, thick] (run) -- (pass); \draw[->, thick] (run) -- (ft);
\node[anchor=north, text=accentc!65!black, font=\scriptsize\bfseries] (fq) at (9.6,-6.15)
  {Are these decisions supported by the behavior of the C program?};
\begin{scope}[on background layer]
  \node[draw=black!40, rounded corners=5pt, fill=black!2, inner sep=8pt, fit=(cmp)(fb)(tests)(pass)(ft)(fq)] (ev) {};
  \node[draw=accentc, very thick, dashed, rounded corners=5pt, fill=accentc!8, inner sep=4pt, fit=(pass)(ft)(fq)] (focus) {};
\end{scope}
\node[anchor=north west, font=\footnotesize\bfseries] at (ev.north west) {Evaluation of each translation};
\end{tikzpicture}
}
\caption{From a C program to a PASS or FAIL decision. The translator is shown as one box, since its internal steps differ between systems and lie outside the scope of this thesis. A translation that compiles is run on the inputs of the benchmark tests and receives PASS when every result matches the expected one. The highlighted decisions are those this thesis re-examines.}
\label{fig:intro-workflow}
\end{figure}

\section{Motivation}\label{sec:intro-motivation}

Nevertheless, generating compilable Rust does not establish that the translation preserves the behavior of its C counterpart. Published evaluations commonly rely on compilation and execution against available test suites, using the resulting PASS/FAIL decisions to assess translation success (Figure~\ref{fig:intro-workflow}). These decisions, however, depend on the inputs exercised and the expectations used to judge them. This raises a question beyond the ability of LLMs to generate Rust code: how reliably do the reported evaluation decisions reflect the behavioral evidence supporting them?

A concrete example of this limitation appears in Reboot, a recent system for translating C-based language interpreters into safe Rust. Wang et al.\ reported that all six translated interpreters passed 100\% of their provided test suites, yet achieved pass rates ranging from 61.58\% to 91.78\% on separately written validation tests that had not been available during translation~\cite{wang2026reboot} (Figure~\ref{fig:intro-signals}a). These results demonstrate that complete success on an existing test suite can coexist with failures exposed by additional tests. Consequently, a reported PASS provides evidence of success under the evaluated inputs, but does not establish that the translated program preserves the behavior of its C counterpart more generally.

\begin{figure}[htbp]
\centering
\resizebox{\linewidth}{!}{\begin{tikzpicture}[>=Stealth, font=\footnotesize]
% ---------- (a) Reboot ----------
\def\W{4.4}
\node[anchor=west, font=\small\bfseries, text=passc] at (-2.0,0.75) {(a) Accepted, yet diverges on held-out tests};
\node[anchor=west, font=\scriptsize] at (-2.0,0.32) {Reboot: six C interpreters translated to Rust};
\foreach \name/\v [count=\i] in {wren/91.78, awk/78.82, gnu-bc/78.57, mujs/74.77, picoc/69.50, pocketpy/61.58}{
  \pgfmathsetmacro{\y}{-0.5*\i}
  \node[anchor=east, font=\scriptsize\ttfamily] at (-0.08,\y) {\name};
  \fill[passc!75] (0,\y-0.17) rectangle ({\W*\v/100},\y+0.17);
  \node[anchor=east, font=\scriptsize, text=white] at ({\W*\v/100},\y) {\v};
}
\draw[black!60] (0,-0.15) -- (0,-3.25) -- (\W,-3.25);
\foreach \t in {0,50,100}{
  \draw[black!60] ({\W*\t/100},-3.25) -- ({\W*\t/100},-3.33)
     node[below, font=\scriptsize] {\t};
}
\draw[dashed, thick, black!75] (\W,-0.15) -- (\W,-3.25);
\node[anchor=south, font=\scriptsize, align=center] at (\W,-0.17) {provided\\tests: 100\%};
\node[font=\scriptsize] at ({\W/2},-4.1) {\% of held-out validation tests passed};
% ---------- (b) Holdem_Odds ----------
\begin{scope}[xshift=8.9cm,
  box/.style={draw, rounded corners=2pt, align=center, font=\scriptsize, minimum width=5.3cm, inner sep=3pt}]
\node[anchor=west, font=\small\bfseries, text=failc] at (-2.7,0.75) {(b) Rejected, yet faithful};
\node[anchor=west, font=\scriptsize] at (-2.7,0.32) {CRUST-Bench, program \texttt{Holdem\_Odds}};
\node[box] (in) at (0,-0.5) {Hand in the test: \texttt{6s 5c 4s 3h 2d}, a straight\\(s, c, h, d: spades, clubs, hearts, diamonds)};
\node[box, fill=black!5] (st) at (0,-1.75) {Test code builds every card\\with \texttt{Suit::Spade}};
\node[box] (hd) at (0,-2.65) {Hand actually tested: $6\spadesuit\ 5\spadesuit\ 4\spadesuit\ 3\spadesuit\ 2\spadesuit$};
\draw[->] (in) -- (st); \draw[->] (st) -- (hd);
\node[anchor=west, font=\scriptsize] at (-2.6,-3.3) {C program};
\node[anchor=west, font=\scriptsize] at (-2.6,-3.7) {Translation};
\node[anchor=west, font=\scriptsize] at (-2.6,-4.1) {Test expects};
\node[anchor=west, font=\scriptsize\ttfamily] at (-1.05,-3.3) {StraightFlush};
\node[anchor=west, font=\scriptsize\ttfamily] at (-1.05,-3.7) {StraightFlush};
\node[anchor=west, font=\scriptsize\ttfamily] at (-1.05,-4.1) {Straight};
\node[anchor=west, font=\scriptsize\itshape, text=black!65] at (1.6,-3.7) {same as C};
\draw[black!40] (-2.6,-4.38) -- (2.6,-4.38);
\node[anchor=west, font=\scriptsize\bfseries] at (-2.6,-4.62) {Verdict};
\node[anchor=west, draw=failc, very thick, text=failc, font=\scriptsize\bfseries, rounded corners=2pt, inner sep=2pt] at (-1.05,-4.62) {FAIL};
\node[anchor=west, font=\scriptsize\itshape, text=black!65] at (-0.1,-4.62) {although the translation matches C};
\node[font=\scriptsize, align=center, text width=5.4cm] at (0,-5.2)
  {The same cause makes 5 of the 24 tests of this program fail (issue \#37).};
\end{scope}
\end{tikzpicture}
}
\caption{Two signals that a test-based decision can be wrong in either
direction. (a)~All six interpreters translated in Reboot pass every
provided test (counted per test file) but only part of a validation suite
that the authors wrote separately, with expected outputs taken from the C
programs, and withheld during translation~\cite{wang2026reboot}.
(b)~A CRUST-Bench test turns a mixed-suit straight into five spades and
still expects \code{Straight}, so a translation that behaves like the C
program would fail, as the analysis in the issue shows~\cite{crustbench-issues}.
Both cases lie outside the corpus studied in this thesis.}
\label{fig:intro-signals}
\end{figure}

The reliability of test-based evaluation can also be affected by the expectations against which translations are judged. An issue reported in CRUST-Bench illustrates this problem through the \texttt{Holdem\_Odds} project, where several tests construct card combinations incorrectly and expect results that conflict with the original C implementation. In one case, a test expects a \texttt{Straight} despite constructing a straight flush, for which the C program returns \texttt{StraightFlush}~\cite{crustbench-issues} (Figure~\ref{fig:intro-signals}b). A Rust translation preserving the C program's behavior would therefore fail the affected test. Although this example does not establish how frequently such problems occur, it demonstrates that a FAIL decision may reflect an inconsistency in the evaluation rather than a behavioral defect in the translation.

Together, these observations raise concerns about the evidential basis of reported translation success rates. Additional tests can challenge decisions to accept translations, while incorrect test expectations can undermine decisions to reject them. However, isolated examples do not establish how frequently either problem affects published evaluations, nor how discrepancies in individual decisions influence the aggregate success rates used to assess translation systems. This motivates a systematic, independent examination of the behavioral evidence underlying reported PASS/FAIL decisions.

\section{Problem Statement and Research Questions}\label{sec:intro-problem}

The evidence presented in Section~\ref{sec:intro-motivation} highlights a fundamental problem in the evaluation of automated C-to-Rust translation: \textbf{published PASS/FAIL decisions may not fully reflect the behavioral agreement between translated Rust programs and their C counterparts}. The problem addressed in this thesis is therefore to determine how strongly these decisions are supported by independently collected behavioral evidence and how reassessing them changes the interpretation of reported translation success.

To investigate this problem, the thesis examines the publicly released translations and evaluation decisions in the HARVEST results repository, using the C programs and test vectors from the TRACTOR Public Test Corpus~\cite{harvest-results,tractor-test-corpus}. The study focuses on the evaluation of LLM-based C-to-Rust translation, while also including the two rule-based translators, C2Rust and Laertes, that the original evaluation reports as baselines on the same corpus. The unit of analysis is the original evaluation decision associated with each case--system pair. The decisions are examined with an independent differential oracle, which compares the behavior of a translation with that of its C program on the same input without reading the expected outputs stored in the tests. The study addresses three research questions:

\begin{description}
\item[RQ1.] Among the translations accepted by their original evaluation, for what proportion does the independent oracle detect a behavioral divergence from the C program?
\item[RQ2.] Among the translations rejected at the test stage by their original evaluation, what proportion of the rejection decisions does the independent oracle find unsupported?
\item[RQ3.] After incorporating the reclassifications supported by the independent oracle, within what interval does each system's success rate over the study population lie, compared with its original rate?
\end{description}

\section{Research Approach}\label{sec:intro-approach}

The thesis answers these questions through a post-hoc re-evaluation. The translations are treated as fixed, so no translator is rerun and no translation is repaired. Accepted translations and translations rejected at the test stage are re-examined, while rejections at the build stage keep their published status.

The re-evaluation proceeds in three steps (Figure~\ref{fig:intro-pipeline}). First, each PASS decision and each FAIL decision at the test stage is reproduced with the public runner of the corpus. Second, a differential testing harness built for this thesis runs the C program and its Rust translation on the same inputs and compares their observable behavior: accepted translations on the existing test vectors and on generated inputs, and test-stage rejections on the vectors that failed. An observation counts as evidence only when it recurs across repeated executions and a sanitizer check of the C execution reports no undefined behavior. Third, the outcomes of both branches are returned to the original population, together with the decisions left without a result, and the success rate of each system after reclassification is stated as an interval. Chapter~\ref{ch:methodology} defines the evidence rules and the metrics.

\begin{figure}[!htbp]
\centering
\resizebox{\linewidth}{!}{\begin{tikzpicture}[>=Stealth, font=\footnotesize,
  st/.style={draw, rounded corners=2pt, align=center, text width=#1, minimum height=12mm, inner sep=3pt},
  st/.default=2.8cm,
  pa/.style={st=#1, draw=passc, fill=passc!8, thick}, pa/.default=2.8cm,
  fa/.style={st=#1, draw=failc, fill=failc!8, thick}, fa/.default=2.8cm,
  outb/.style={draw, thick, rounded corners=2pt, align=center, text width=#1, minimum height=12mm},
  outb/.default=1.1cm]
\node[st=2.3cm, fill=black!4] (src) at (-0.35,0) {HARVEST and TRACTOR: released translations, decisions, C programs, test vectors};
\node[st=1.9cm] (rep) at (2.45,0) {Reproduce the decisions with the TRACTOR runner};
\node[pa] (a1) at (6.3,1.55) {Run the C program and the translation on test vectors and generated inputs};
\node[pa] (a2) at (9.55,1.55) {Retain divergences that recur and pass the sanitizer check};
\node[outb, draw=passc] (a3) at (12.3,1.55) {RQ1};
\node[fa] (b1) at (6.3,-1.55) {Replay the failing vectors on the C program and the translation};
\node[fa] (b2) at (9.55,-1.55) {Adjudicate each vector: repeated runs, sanitizers};
\node[outb, draw=failc] (b3) at (12.3,-1.55) {RQ2};
\node[outb=1.9cm] (r3) at (14.6,0) {RQ3:\\success-rate interval};
\draw[->] (src) -- (rep);
\draw[->, passc, thick] (rep.east) -- ++(0.3,0) |- node[pos=0.78, above, font=\scriptsize\bfseries]{PASS} (a1.west);
\draw[->, failc, thick] (rep.east) -- ++(0.3,0) |- node[pos=0.78, below, font=\scriptsize\bfseries]{FAIL} (b1.west);
\draw[->, passc, thick] (a1) -- (a2); \draw[->, passc, thick] (a2) -- (a3);
\draw[->, failc, thick] (b1) -- (b2); \draw[->, failc, thick] (b2) -- (b3);
\draw[->] (a3.east) -- (r3); \draw[->] (b3.east) -- (r3);
\draw[->, dashed, black!60] (rep.south) -- (rep.south |- 0,-2.75) -- node[pos=0.5, below, font=\scriptsize]{decisions left open widen the RQ3 interval} (r3.south |- 0,-2.75) -- (r3.south);
\end{tikzpicture}
}
\caption{Overview of the re-evaluation in three steps. Each PASS decision and each test-stage FAIL decision is first reproduced. Accepted translations and test-stage rejections then follow separate branches with different inputs, and RQ3 combines the results of both branches with the decisions left without a result.}
\label{fig:intro-pipeline}
\end{figure}

\section{Research Contributions}\label{sec:intro-contributions}

This thesis makes a scientific and a technical contribution to the evaluation of automated C-to-Rust translation, with a focus on the validity of test-based evaluation decisions, and its findings have practical significance.

\paragraph{Scientific contribution.} The study contributes a re-evaluation procedure for published PASS/FAIL decisions and the empirical evidence obtained with it. The procedure states, for acceptances and for test-stage rejections, when the observed behavior of a translation and its C program counts as evidence for or against a decision, and it computes the success rate after reclassification as an interval over the original population, so that the decisions that could not be examined remain visible in the reported rate (Chapter~\ref{ch:methodology}). Applied to the released translations of ten systems, the procedure provides empirical evidence on the validity of published test-based PASS/FAIL decisions in C-to-Rust translation. By independently examining both accepted translations and test-stage rejections, it quantifies how often these decisions are contradicted or left unsupported by behavioral evidence and what this implies for reported success rates (Chapter~\ref{ch:results}). The analysis also accounts for unresolved evaluations, distinguishing observed discrepancies from conclusions that cannot be established with the available evidence.

\paragraph{Technical contribution.} The thesis develops a differential testing harness that implements this procedure for executable programs and library functions (Section~\ref{sec:m-harness}). The harness integrates input generation, paired C--Rust execution, repeated observations, sanitizer checks, and evidence recording, and its correctness was checked with held-out decisions, mutants and unit tests (Section~\ref{sec:m-evidence}). Together with the classification and accounting rules (Appendix~\ref{app:classes}), it makes the re-evaluation repeatable, and its code and result files are publicly available (Appendix~\ref{app:s-env}).

\paragraph{Practical significance.} The findings inform how automated code translation systems can be evaluated and how their results should be reported. They indicate the value of testing accepted translations beyond existing test suites and of reporting evaluation coverage and unresolved cases explicitly (Section~\ref{sec:d-implications}). These practices can help researchers and developers interpret reported translation success rates more critically and improve the transparency of future evaluations. The findings serve the interpretation of the evaluation results of the systems studied and do not imply a ranking of those systems.

\section{Thesis Organization}\label{sec:intro-organization}

Table~\ref{tab:intro-organization} summarizes the content of the remaining chapters.

\begin{table}[!htb]
\centering
\small
\caption{Organization of the thesis.}
\label{tab:intro-organization}
\begin{tabularx}{\linewidth}{@{}p{0.40\linewidth}X@{}}
\toprule
Chapter & Content \\
\midrule
\textbf{Chapter 2: Related Work} & Evaluation approaches for C-to-Rust translation, differential testing, studies that re-evaluate published results, and the research gap. \\
\textbf{Chapter 3: Methodology} & Design of the re-evaluation, study population, the harness, input generation and execution, evidence rules, and metrics. \\
\textbf{Chapter 4: Results} & Findings for RQ1, RQ2 and RQ3, with the origin and weight of the divergences, the decisions without a result, and the effect of the comparison rule. \\
\textbf{Chapter 5: Discussion} & Interpretation of the findings, relation to earlier studies, threats to validity, and implications for the evaluation of translations. \\
\textbf{Chapter 6: Conclusion} & Answers to the research questions and directions for future work. \\
\bottomrule
\end{tabularx}
\end{table}
